\documentclass[11pt,a4paper]{article}
\usepackage[margin=20mm]{geometry}
\usepackage{amsmath,amssymb,mathtools}
\usepackage{enumitem}
\usepackage{xcolor}
\usepackage{fontspec}
\setmainfont{Arial Unicode MS}
\setmonofont{Menlo}
\definecolor{navy}{HTML}{172554}
\definecolor{blue}{HTML}{1D4ED8}
\setlength{\parskip}{0.65em}
\setlength{\parindent}{0pt}
\newcommand{\heading}[1]{\section*{\color{blue}#1}}
\newcommand{\question}[1]{\subsection*{#1}}
\newcommand{\code}[1]{\texttt{#1}}
\begin{document}

\begin{center}
{\Large\color{navy} Transformer From Scratch Lab}\par
\vspace{8mm}
{\huge\color{navy} Week 0 Solutions}\par
\vspace{4mm}
{\Large 先看懂信息怎样流动}
\end{center}

\textbf{使用说明：}只在完成 Week 0 Problem Set、Shape Checker Lab 和自己的反馈记录后阅读。\\
本文件展示一种完整推导与工程解释；它不是唯一可接受的表述方式。

\heading{Part A：概念与推导}
\question{A1. Token 表与 shape}
因为
\[
X \in \mathbb{R}^{3 \times 2}, \qquad W \in \mathbb{R}^{2 \times 4},
\]
内维度 $2$ 相等，所以
\[
XW \in \mathbb{R}^{3 \times 4}.
\]
输出的第 $i$ 行仍对应输入的第 $i$ 个 token；第 $j$ 列是权重矩阵第 $j$ 列定义的一种新特征。若 $X_{2,1}$ 增加 $1$，只会改变输出第 $2$ 行的元素，因为矩阵乘法的第 $i$ 行只由 $X$ 的第 $i$ 行参与计算。

\question{A2. 矩阵乘法}
\[
A = \begin{bmatrix} 1 & 2 \\ 0 & 3 \\ 4 & 1 \end{bmatrix},
\qquad
B = \begin{bmatrix} 2 & 1 \\ 1 & 0 \end{bmatrix}.
\]
第 $1$ 行为
\[
[1,2]B = [1\cdot2+2\cdot1,\;1\cdot1+2\cdot0] = [4,1].
\]
第 $3$ 行为
\[
[4,1]B = [4\cdot2+1\cdot1,\;4\cdot1+1\cdot0] = [9,4].
\]
一般地，若 $A \in \mathbb{R}^{m \times n}$、$B \in \mathbb{R}^{n \times p}$，则
\[
(AB)_{i,j}=\sum_{k=1}^{n} A_{i,k}B_{k,j}.
\]
同一个 $B$ 作用于 $A$ 的每一行，因此每个 token 都经过相同参数变换；这正是后续 $XW_Q$、$XW_K$ 和 $XW_V$ 的形式。

\question{A3. Dot product}
\[
u\cdot v = 1\cdot3+2\cdot(-1)=1,
\qquad
u\cdot w = 1\cdot1+2\cdot1=3.
\]
在这个例子中，$w$ 比 $v$ 更匹配 $u$。但 $1$ 和 $3$ 只是未归一化 score：它们没有保证非负、没有保证总和为 $1$，也无法直接说明应该读取多少 value。

\question{A4. Stable softmax}
对 $z=[2,1,0]$，
\[
\operatorname{softmax}(z)=\frac{1}{e^2+e+1}[e^2,e,1].
\]
三项之和为 $1$。对任意 $c$，
\[
\operatorname{softmax}(z+c\mathbf{1})_i
=\frac{e^{z_i+c}}{\sum_j e^{z_j+c}}
=\frac{e^c e^{z_i}}{e^c\sum_j e^{z_j}}
=\operatorname{softmax}(z)_i.
\]
因此可令 $c=-\max_i z_i$。这样最大 exponent 为 $e^0=1$，避免 $e^{1000}$ 一类溢出，同时不改变输出概率。

\heading{Part B：使用边界与全局位置}
\question{B1. 不合法的 shape}
例如
\[
A\in\mathbb{R}^{3\times2}, \qquad B\in\mathbb{R}^{4\times1}.
\]
左矩阵列数为 $2$，右矩阵行数为 $4$，所以 $AB$ 未定义。Shape Checker 应报告 `shape/维度` 类别，指出 left 的列数 $2$ 与 right 的行数 $4$ 不匹配。原始异常可能告诉我们运算失败，却不一定把它关联为“内维度契约被破坏”；后者才是能迁移到 Q/K/V 投影的机制信息。

\question{B2. Softmax 的错误轴}
若 $S_{i,j}$ 是 query $i$ 对 key $j$ 的 score，沿最后一维归一化给出
\[
P_{i,j}=\frac{e^{S_{i,j}}}{\sum_k e^{S_{i,k}}},
\]
即同一个 query 在所有 key 间分配读取权重。沿第 $0$ 维归一化则让不同 query 对同一 key 竞争，这不再表示每个 query 自己的读取分配。

例如
\[
S=\begin{bmatrix}4&0\\0&4\end{bmatrix}.
\]
沿行归一化时，每个 query 更偏向自己的对应 key；沿列归一化时，每个 key 在不同 query 间分配权重，语义已改变。

\question{B3. 全局连接}
完整链可分为两步：
\[
\begin{aligned}
X &\xrightarrow{\text{projection}} (Q,K,V)
\xrightarrow{\text{dot products}} QK^\top,\\
QK^\top &\xrightarrow{\text{row-wise softmax}}
\operatorname{softmax}\!\left(\frac{QK^\top}{\sqrt{d_k}}\right)
\xrightarrow{\text{weighted sum}}
\operatorname{softmax}\!\left(\frac{QK^\top}{\sqrt{d_k}}\right)V.
\end{aligned}
\]
Week 0 已经覆盖 token 表的行/列语义、投影的矩阵乘法、score 的 dot product、将 score 变为权重的 softmax，以及输出 shape 的预测。Week 1 的新增工作不是引入完全陌生的算子，而是将这些操作按 attention 的顺序组合，并验证每一行权重如何决定对 $V$ 的加权读取。

\heading{Part C：Lab 后工程问题}
\question{C1. 函数契约}
\code{validate\_matrix\_shape} 的单一职责是验证一个 shape 是否二维、每个维度是否为正，\\
并在错误消息中保留 left/right 来源。例如 $(3,)$ 应在这里被拒绝。\\
\code{matrix\_product\_shape} 的单一职责是先验证两侧，再检查内维度并在合法时返回 $(m,q)$。\\
例如 $(3,2)$ 与 $(4,1)$ 的内维度错误属于后者。拆分后，错误分类可测试且可解释。

\question{C2. 诊断样例}
一个可复现的主动错误是输入 left shape $(3,2)$ 与 right shape $(4,1)$。\\
现象是矩阵乘法不能定义；假设是内维度不相等。\\
最小测试应断言函数抛出 \code{ValueError}，并包含 left、right 以及数值 $2$、$4$。\\
修复不是改变矩阵数值，而是让右矩阵改为 $(2,1)$ 或调整左矩阵列数。\\
这个例子说明 shape 是接口契约，不是调试时可随意修改的标签。

\question{C3. 迁移到 Week 1}
对
\[
Q=XW_Q,\qquad K=XW_K,\qquad V=XW_V,
\]
Shape Checker 可在每次投影前确认 $X$ 的最后一维等于 $W_Q$、$W_K$、$W_V$ 的首维。\\
它不能防止的例子是：三个权重矩阵 shape 全部合法，却错误地令 $W_Q=W_K=W_V$，\\
或对 $QK^\top$ 沿错误轴做 softmax。此时 shape 正确，但 attention 机制仍可能错误。

\vfill
\begin{center}\small Week 0 Solutions\end{center}
\end{document}
